% Part: incompleteness
% Chapter: representability-in-q
% Section: composition-representable

\documentclass[../../../include/open-logic-section]{subfiles}

\begin{document}

\olfileid{inc}{req}{cmp}
\olsection{फलनांचे संयोजन~$\Th{Q}$ मध्ये निरूपणीय असते}

समजा $h$ पुढीलप्रमाणे परिभाषित आहे:
\[
h(x_0,\dots,x_{l-1}) = f(g_0(x_0,\dots,x_{l-1}), \dots,
g_{k-1}(x_0,\dots,x_{l-1})).
\]
येथे $f$ आणि अनुक्रमे $g_0$,
\dots,~$g_{k-1}$ या फलनांचे निरूपण करणारी
!!{formula}s $!A_f, !A_{g_0}, \dots,
!A_{g_{k-1}}$ आपल्याला आधीच सापडली आहेत. आता $h$ चे
निरूपण करणारे !!a{formula}~$!A_h$ शोधायचे आहे.

सर्व फलनांचे $1$ आदान असलेली सोपी बाब आधी पाहू; म्हणजे
$h(x) = f(g(x))$ विचारात घेऊ. $!A_f(y, z)$ हे~$f$ चे,
आणि $!A_g(x, y)$ हे~$g$ चे निरूपण करत असेल, तर $h$ चे
निरूपण करणारे !!a{formula}~$!A_h(x, z)$ शोधायचे आहे.
$h(x) = z$ तेव्हाच आणि तरच असा एखादा~$y$ असतो की
$z = f(y)$ आणि $y = g(x)$ ही दोन्ही विधाने खरी असतात.
($h(x) = z$ असेल तर $g(x)$ हा असा~$y$ असतो; आणि असा
$y$ असेल, तर $y = g(x)$ आणि $z = f(y)$ असल्याने
$z = f(g(x))$.) यावरून $!A_h(x, z)$ साठी
$\lexists[y][(!A_g(x, y) \land !A_f(y,
  z))]$ हे योग्य सूत्र असावे असे दिसते. आता फक्त
$\Th{Q}$ संबंधित !!{formula}s सिद्ध करते, हे तपासायचे आहे.

\begin{prop}
\ollabel{prop:rep1}
$h(n) = m$ असेल, तर $\Th{Q} \Proves !A_h(\num{n}, \num{m})$.
\end{prop}

\begin{proof}
$h(n) = m$, म्हणजे $f(g(n)) = m$, असे समजा. $k = g(n)$ असू द्या.
मग
\begin{align*}
  \Th{Q} & \Proves !A_g(\num{n}, \num{k})
  \intertext{कारण $!A_g$ हे~$g$ चे निरूपण करते; तसेच}
  \Th{Q} & \Proves !A_f(\num{k}, \num{m})
  \intertext{कारण $!A_f$ हे~$f$ चे निरूपण करते. म्हणून}
  \Th{Q} & \Proves !A_g(\num{n}, \num{k}) \land !A_f(\num{k}, \num{m})
  \intertext{आणि त्यावरून}
  \Th{Q} & \Proves \lexists[y][(!A_g(\num{n}, y) \land !A_f(y, \num{m}))],
\end{align*}
म्हणजेच $\Th{Q} \Proves !A_h(\num{n}, \num{m})$.
\end{proof}

\begin{prop}
\ollabel{prop:rep2}
$h(n) = m$ असेल, तर $\Th{Q} \Proves \lforall[z][(!A_h(\num{n}, z) \lif
  z = \num{m})]$.
\end{prop}

\begin{proof}
$h(n) = m$, म्हणजे $f(g(n)) = m$, असे समजा. $k = g(n)$ असू द्या.
मग
\begin{align*}
  \Th{Q} & \Proves \lforall[y][(!A_g(\num{n}, y) \lif \eq[y][\num{k}])]
  \intertext{कारण $!A_g$ हे~$g$ चे निरूपण करते; तसेच}
  \Th{Q} & \Proves \lforall[z][(!A_f(\num{k}, z) \lif \eq[z][\num{m}])]
  \intertext{कारण $!A_f$ हे~$f$ चे निरूपण करते. थोडेसे
    तर्कशास्त्र वापरून पुढील विधानही दाखवता येते:}
  \Th{Q} & \Proves \lforall[z][(\lexists[y][(!A_g(\num{n}, y) \land
      !A_f(y, z))] \lif \eq[z][\num{m}])].
\end{align*}
म्हणजेच $\Th{Q} \Proves \lforall[y][(!A_h(\num n, y) \lif \eq[y][\num m])]$.
\end{proof}

$f$ आणि~$g_i$ यांची आदाने $1$ पेक्षा अधिक असलेल्या गुंतागुंतीच्या
बाबतीतही हीच कल्पना लागू पडते.

\begin{prop}
\ollabel{prop:rep-composition}
$!A_f(y_0, \dots, y_{k-1}, z)$ हे~$\Th{Q}$ मध्ये
$f(y_0, \dots, y_{k-1})$ चे, आणि
$!A_{g_i}(x_0, \dots, x_{l-1}, y)$ हे~$\Th{Q}$ मध्ये
$g_i(x_0, \dots, x_{l-1})$ चे निरूपण करत असेल, तर
\begin{multline*}
  \lexists[y_0\dots][\lexists[y_{k-1}][(!A_{g_0}(x_0,\dots,x_{l-1},y_0) \land
      \dots \land {}]]\\
  !A_{g_{k-1}}(x_0,\dots,x_{l-1},y_{k-1}) \land !A_f(y_0,\dots,y_{k-1},z))
\end{multline*}
हे सूत्र पुढील फलनाचे निरूपण करते:
\[
h(x_0, \dots, x_{l-1}) = f(g_0(x_0, \dots, x_{l-1}), \dots, g_{k-1}(x_0,
\dots, x_{l-1})).
\]
\end{prop}

\begin{proof}
सरावासाठी.
\end{proof}

\begin{prob}
\olref[inc][req][cmp]{prop:rep1} आणि
\olref[inc][req][cmp]{prop:rep2} यांच्या पुराव्यांचा आधार घेऊन
\olref[inc][req][cmp]{prop:rep-composition} चा पुरावा तपशीलवार द्या.
\end{prob}

\end{document}
